%请使用pdflatex编译 
%代码是我自己写的，但是方法啥的全是网上搜的，我也不知道我这种情况有没有版权QAQ，转载标明出处即可
%联系作者：邮箱2978028549@qq.com，B站ID“我不是wc"
%文档首发于B站
%求打赏
\documentclass{article}
\usepackage{ctex}                                           %显示中文，更改字体
\usepackage{setspace}                                       %调整页行距
\usepackage{geometry}                                       %设置页边距
\usepackage{booktabs}                                       %制作三线表
\usepackage{graphicx}                                       %插入图片
\usepackage[CJKbookmarks=true,pdfpagemode=UseOutlines,bookmarksnumbered,bookmarksopen]{hyperref}                    %生成书签
\hypersetup{hidelinks,colorlinks=true,allcolors=black,pdfstartview=Fit,breaklinks=true}%配置书签格式
\title{文章的标题}                                          %文章标题
\author {作者名称}                                          %作者的名称
\date{\today}                                               %日期article
\renewcommand{\baselinestretch}{1.2}                        %调整行距为1.2倍标准行距
\renewcommand\refname{\ }                                   %吧"reference"改成"参考文献"
\renewcommand{\abstractname}{摘要}                          %把“abstract”改成“摘要”
%\setlength{\baselineskip}{20pt}                            %调整行距为20pt,(不好用，如果需要使用的话要放在需要修改的文字前面)
%\geometry{left=2.5cm,right=2.5cm,top=2.5cm,bottom=2.5cm}   %设置页边距

\begin{document}
\setlength{\parindent}{0em}
\maketitle                                                  %生成文章题目，作者，时间
\thispagestyle{empty}                                       %本页不生成页码
\begin{abstract}                                            %摘要 开始
在这里写摘要\\ 摘要\\ 摘要
\end{abstract}                                              %摘要 结束
\newpage                                                    %换页
\tableofcontents                                            %生成目录
\thispagestyle{empty}                                       %本页不生成页码
\newpage                                                    %换页
\setcounter{page}{1}                                        %文章页码从1开始重新编排
\section{问题重述}                                          %换页
\newpage
\section{问题分析}
\newpage                                                    %换页
\section{模型假设与约定}
\newpage                                                    %换页
\section{符号说明与名词定义}
\newpage                                                    %换页
\section{模型建立}
\newpage                                                    %换页
\section{模型求解}
\newpage                                                    %换页
\section{模型建议}
\newpage                                                    %换页
\section{模型评价}
\newpage                                                    %换页
\section{模型推广}
\newpage
\section{参考文献1}
\cite{ref1}引用参考文献1                                     %换页
\section{参考文献2}
\renewcommand{\baselinestretch}{0.01}                       %调整行距为1.2倍标准行距
\begin{thebibliography}{99}                                 %最大摘要上限为99个
\bibitem{ref1}这是第一个参考文献
\bibitem{ref2}这是第二个参考文献
\bibitem{ref3}这是第三个参考文献
\bibitem{ref4}这是第四个参考文献
\end{thebibliography}
\newpage                                                    %换页
\section{附录}
\section{首行如何空两格}
%由于前面使用了代码\setlength{\parindent}{0em}，首行不再缩进，如：
up主是完美男人up主是完美男人up主是完美男人up主是完美男人up主是完美男人up主是完美男人up主是完美男人up主是完美男人up主是完美男人up主是完美男人up主是完美男人up主是完美男人up主是完美男人up主是完美男人\\
************************************************************************\\
%我们可以利用\indent\qquad\ 实现空两格，如:
\indent\qquad\  up主是完美男人up主是完美男人up主是完美男人up主是完美男人up主是完美男人up主是完美男人up主是完美男人up主是完美男人up主是完美男人up 主是完美男人\\
\newpage                                                    %换页
\section{如何更改字体}
%局部更改字体，例：
{\songti    宋体}局部字体\\
{\heiti     黑体}局部字体\\
{\fangsong  仿宋}局部字体\\
{\kaishu    楷书}局部字体\\
************************************************************************\\
%全局更改宋体，例：
\songti    宋体  全局字体\\
\heiti     黑体  全局字体\\
\fangsong  仿宋  全局字体\\
\kaishu    楷书  全局字体\\
************************************************************************\\
%字号同理
\zihao{3}   三号字体\\
\zihao{1}   一号字体\\
{\zihao{3}  局部三号} 还是一号\\
\newpage
\section{制作三线表}
\zihao{5}
\begin{center}
\begin{tabular}{ccc}
    \toprule
    B站搜索我不是wc&B站搜索我不是wc&B站搜索我不是wc\\
    \midrule
    B站搜索我不是wc&B站搜索我不是wc&B站搜索我不是wc\\
    B站搜索我不是wc&B站搜索我不是wc&B站搜索我不是wc\\
    B站搜索我不是wc&B站搜索我不是wc&B站搜索我不是wc\\
    B站搜索我不是wc&B站搜索我不是wc&B站搜索我不是wc\\
    B站搜索我不是wc&B站搜索我不是wc&B站搜索我不是wc\\
    \bottomrule
\end{tabular}
\end{center}
\newpage
\section{插入图片}
\begin{figure}[hbt]%先尝试能否将照片放在原处(here)，如果不能，那么放在页面底部(bottom),最后尝试放在顶部(top)
  \centering
  \includegraphics[scale=0.2]{1.jpg}%将照片1.jpg，缩放成原来的0.2倍
  \caption{标题}
\end{figure}  
\newpage
\section{并排插入两个图片}
\begin{figure}[hbt]
	\begin{minipage}[t]{0.5\textwidth}
		\centering
		\includegraphics[scale=0.2]{1.jpg}
		\caption{标题1}
	\end{minipage}
	\qquad
	\begin{minipage}[t]{0.5\textwidth}
		\centering
		\includegraphics[scale=0.2]{1.jpg}
		\caption{标题2}
	\end{minipage}
\end{figure}
\newpage
\section{书签乱码解决办法}
\begin{figure}[hbt]%先尝试能否将照片放在原处(here)，如果不能，那么放在页面底部(bottom),最后尝试放在顶部(top)
  \centering
  \includegraphics[scale=0.5]{2.png}%将照片2.png，缩放成原来的0.2倍
  \caption{标题}
\end{figure} 
该方法来自https://blog.csdn.net/he\_nan/article/details/106165716
\newpage
\end{document}
